% ===========================================================================
%  sviluppatore/07-marketplace/01-identita.tex
% ===========================================================================
\section{Il ponte verso il sito}\label{sez:ponte}

La dashboard non sa chi è l'utente: lo chiede al sito. Non conserva password
né segreti: li tiene il sito. Tutto ciò che riguarda identità, denaro e
catalogo passa da \file{dashboard/app\_api.py}, e il meccanismo è un token
personale.

\subsection{Dove sta il token}

\begin{codiceInterno}{python}{La ricerca del token, in ordine}{dashboard/app\_api.py:67}
def token() -> str | None:
    global _TOKEN, _TOKEN_LOADED
    if _TOKEN_LOADED:
        return _TOKEN
    _TOKEN_LOADED = True
    env = (os.environ.get("LYRA_APP_TOKEN") or "").strip()
    if env:
        _TOKEN = env
        return _TOKEN
    try:
        p = _token_path()
        if p.exists():
            data = json.loads(p.read_text(encoding="utf-8"))
            _TOKEN = (data.get("token") or "").strip() or None
    except Exception:
        _TOKEN = None
    return _TOKEN
\end{codiceInterno}

\noindent
Variabile d'ambiente prima, file dopo, e il risultato è messo in cache per
tutta la vita del processo. La cache è la ragione per cui, dopo un accesso
nuovo, può servire un riavvio: \cod{\_TOKEN\_LOADED} resta vero.

\subsection{La chiamata che non solleva}

\begin{commento}{dashboard/app\_api.py:137 $\cdot$ \cod{chiama\_sito()}}
Chiamata al sito che NON solleva: ritorna (ok, dati).

Serve alle rotte che fanno solo da ponte — sicurezza dell'account, per
esempio. L'app non conserva password ne' segreti: li tiene il sito, e qui
si inoltra soltanto aggiungendo il token di chi ha fatto l'accesso.

Gli errori del sito arrivano con il proprio messaggio e il proprio codice
di stato: tradurli qui significherebbe perdere spiegazioni che il sito ha
scritto apposta.
\end{commento}

\begin{perche}{Perché gli errori del sito si inoltrano senza tradurli}
Il sito scrive messaggi pensati per essere letti dall'utente: «il tuo
abbonamento è scaduto il \ldots», «questa email è già registrata». Sostituirli
con un generico «errore del server» cancella l'informazione. La dashboard fa
da ponte, e un ponte non riscrive ciò che trasporta.
\end{perche}

\subsection{La regola del token giusto}

C'è un secondo segreto, oltre al token dell'utente: quello della dashboard.
Servono a cose diverse, e scambiarli produce il più insidioso dei difetti
di questo capitolo.

\begin{commento}{dashboard/community.py:76 $\cdot$ \cod{SI()}}
SI CHIAMA CON IL TOKEN DELL'UTENTE, non col segreto della dashboard.
Prima no, e per questo non ha mai funzionato: \percorsoHTTP{/api/app/abbonamento}
riconosce la persona dal token, quindi rispondeva 401, l'eccezione veniva
inghiottita e il risultato era «non abbonato» per tutti. Il sintomo era
invisibile perche' in beta gli unici a pubblicare erano gli
amministratori, che passano da un'altra strada.
\end{commento}

\begin{attenzione}{Tre cose insieme rendono un difetto invisibile}
Una risposta \cod{401}, un'eccezione inghiottita, e un ripiego prudente che
risponde «non abbonato». Presi uno per uno sono ragionevoli; insieme
producono un sistema che funziona male e non lo dice. Chi aggiunge una
chiamata al sito controlli quale dei due segreti serve, e verifichi il
percorso d'errore con l'utente giusto --- non con un amministratore, che
spesso passa da un'altra strada e non lo attraversa mai.
\end{attenzione}

\subsection{«Non sapere vale come no»}

Per la pubblicazione il ripiego prudente è invece la scelta corretta, ed è
dichiarata:

\begin{commento}{dashboard/community.py:72 $\cdot$ \cod{In()}}
In caso di errore risponde «livello 0, non puo' pubblicare»: non sapere
vale come no. Pubblicare per errore mette in vetrina roba di cui nessuno
risponde e si toglie dopo; non pubblicare si risolve riprovando.
\end{commento}

\noindent
Il criterio è l'asimmetria del danno: una pubblicazione sbagliata si corregge
con fatica e intanto è pubblica; una pubblicazione mancata si corregge
riprovando. Si noti che è il contrario della regola vista per le \emph{letture}
(\vediSez{pagina-nuova}): lì «non so» non va mostrato come «no», perché
nessuna azione irreversibile ne dipende. La differenza fra i due casi è se
c'è un'azione, non se c'è un'incertezza.
